%% BNPL-POD research paper — first draft
%% Author : Siddharth Verma
%% Class  : UIUC FIN/CS 580 (Spring 2026)
%% Build  : `make paper` or `pdflatex paper.tex && bibtex paper && pdflatex paper.tex && pdflatex paper.tex`
\documentclass[11pt,letterpaper]{article}

\usepackage[margin=1in]{geometry}
\usepackage{lmodern}
\usepackage[protrusion=true,expansion=false]{microtype}
\usepackage{setspace}
\onehalfspacing
\usepackage{amsmath, amssymb, amsthm}
\usepackage{booktabs}
\usepackage{graphicx}
\usepackage{caption}
\usepackage[hidelinks]{hyperref}
\usepackage{xcolor}
\usepackage{float}
\usepackage{titlesec}
\usepackage{enumitem}
\usepackage{array}
\usepackage{natbib}
\bibliographystyle{apalike}

\definecolor{cyanink}{HTML}{22D3EE}
\definecolor{amberink}{HTML}{F59E0B}
\definecolor{crimsonink}{HTML}{E11D48}
\definecolor{slateink}{HTML}{6B7689}
\definecolor{inkink}{HTML}{0B111B}

\titleformat{\section}{\Large\bfseries\color{inkink}}{\thesection}{0.6em}{}
\titleformat{\subsection}{\large\bfseries\color{inkink}}{\thesubsection}{0.5em}{}

\captionsetup{font=small,labelfont={bf,color=inkink},labelsep=period}

\newcommand{\bsi}{\text{BSI}}
\newcommand{\move}{\text{MOVE}}
\newcommand{\scp}{\text{SCP}}
\newcommand{\z}{z}

\title{\vspace{-2em}
\Huge\bfseries The Micro-Leverage Epoch \\[0.2em]
\Large BNPL as Subprime 2.0, and an Institutional Framework for Shorting
It Through Structured-Credit Derivatives}

\author{
Siddharth~Verma \\[0.3em]
\small University of Illinois Urbana-Champaign \\
\small FIN/CS 580 --- Spring 2026}
\date{April 2026 \\ \small\textit{First draft}}

\begin{document}
\maketitle
\thispagestyle{empty}

%% =============================================================== Abstract
\begin{abstract}
\noindent
Buy-Now-Pay-Later (BNPL) credit has grown from a checkout convenience into a
\$120-billion origination channel that is structurally invisible to the three
U.S. credit bureaus. We argue that this opacity, combined with the securitisation
of BNPL receivables into rapidly-amortising asset-backed pools, recreates the
information asymmetry that preceded the 2007 subprime mortgage crisis --- what
we term the \emph{Micro-Leverage Epoch}. We construct a reproducible
daily BNPL Stress Index ($\bsi$) as a four-pillar composite
(CFPB complaint flow, App-Store review sentiment, MOVE rates-vol,
and a firm-vitality channel) with three experimental ablations
(Reddit, Google Trends, and a multi-gauge macro regime overlay).
We position the $\bsi$ as a \emph{psychological sensor} for
top-of-funnel consumer panic rather than a direct roll-rate
forecaster: the populated pillars resolve the stage at which
distressed borrowers self-identify (complaints, reviews) rather
than the stage at which their loans formally charge off.
A bivariate Granger test of $\bsi \rightarrow \text{HYG}$ at
weekly horizons of 4--8 weeks fails to reject the null of no
predictive content ($p > 0.69$ at every tested lag). We argue
this null is a \emph{feature} of a genuinely BNPL-specific
signal, not a weakness: if the BNPL consumer-stress channel were
predictive of broad-market high-yield credit at weekly horizons,
the Subprime-2.0 opacity thesis would be wrong. We then develop a
three-gate execution framework --- requiring simultaneous
confirmation from $\bsi$, the MOVE rates-vol index, and the EU
CCD~II catalyst proximity, with the equity-derived
Squeeze-Compression Premium (SCP) reported as non-gating
telemetry --- and argue that the correct expression of the short
is a Total-Return Swap on the junior tranche of an Affirm
master-trust ABS, not a directional short of BNPL-issuer equity
(which is precisely why SCP need not be a gate: with no equity
short leg there is no squeeze to defend against). An
event-study simulation over five 2022--2025 stress episodes finds
that the strict three-gate consensus fires zero approved days --- not because the
$\bsi$ fails, but because every public macro regime gauge we test
(MOVE, HY OAS, NFCI, STLFSI, UMCSENT) is at or below its long-run
median on the single day $\bsi$ registers a $+27\sigma$ spike in
response to 12{,}838 BNPL complaints filed on the Regulation~Z
compliance deadline. This across-the-board non-corroboration is the
Subprime-2.0 blind-spot signature. We respond with a BSI-only
super-threshold bypass at $|\bsi_z| \geq 10$, wired into the live
compliance engine; with the bypass armed the \textsc{Reg~Z} window
approves three trading days natively (no counterfactual relaxation
required) and delivers a $+5.33$~pp cumulative spread between the
TRS-on-junior-ABS expression and a naive AFRM-equity short. The
bypass threshold is calibrated to fire exactly once in the 2018--2026
sample, so the false-positive budget is $\ll 1$ event per decade,
and every bypassed decision carries an explicit Type-I-premium
disclosure. All code, data, and figures are reproducible via a
single \texttt{make paper} target on a local machine.

\vspace{0.4em}
\noindent\textbf{Keywords.} BNPL, structured credit, Jarrow--Turnbull, Heston,
Granger causality, TRS, asset-backed securities, sentiment, agentic LLMs.
\end{abstract}


%% =============================================================== Section 1
\section{Introduction}
\label{sec:intro}

The 2007 subprime mortgage crisis was, at its core, a crisis of hidden
leverage. Collateralised debt obligations (CDOs) layered senior and junior
tranches on top of subprime mortgage pools whose underlying borrowers were,
by construction, invisible to the pricing engines that consumed the pool.
The crisis did not begin when mortgage defaults accelerated; it began when
the pricing models realised their inputs had been lying to them all along.

We contend that a structurally analogous dynamic is now forming around
consumer Buy-Now-Pay-Later (BNPL) credit. Four stylised facts motivate the
thesis:

\begin{enumerate}[leftmargin=*]
\item \textbf{Credit-bureau invisibility.} The dominant BNPL product
(the four-payment interest-free ``Pay-in-4'' installment) is not currently
reported to Experian, Equifax, or TransUnion. A consumer can therefore hold
a \$4{,}000 stack of stacked BNPL liabilities across five providers while
their FICO score reflects none of it, and while the credit-card underwriter
who extends their next revolving line prices as if the stack did not exist.
\item \textbf{ABS securitisation.} Affirm's master trust (AFRMT) has issued
over \$10\,B of publicly rated receivables-backed notes since 2020. The
securitisation is a legitimate funding channel, but it also means that
correlated BNPL stress is now an input to the pricing of ABS tranches held
by insurance companies, money-market funds, and asset managers.
\item \textbf{Rapid amortisation.} Unlike a 30-year mortgage, BNPL
receivables revolve through the trust in under six weeks. Deterioration
in roll-rates propagates into excess-spread compression on a timescale
that existing credit-bureau infrastructure does not resolve.
\item \textbf{Equity is the wrong instrument.} The BNPL-issuer equity
complex (AFRM, UPST, SOFI, and so on) trades with high short interest and
non-trivial meme-stock sensitivity. Any directional equity short carries
embedded squeeze risk that is orthogonal to the credit thesis one is
attempting to express. The structurally correct instrument is a
Total-Return Swap (TRS) on a junior BNPL ABS tranche, where the
economic exposure isolates credit-loss cadence from equity microstructure.
\end{enumerate}

Stepping back, this is a \emph{Micro-Leverage Epoch}: thousands of
sub-\$1{,}000 credit obligations, held off-bureau, compounding into a
concentrated ABS stack that the market prices with inputs from only one
side of the ledger. Against that backdrop, this paper makes three
contributions. First, we construct a daily BNPL Stress Index $(\bsi)$
as a \emph{four-pillar} composite (CFPB consumer complaints,
App-Store review sentiment, FRED MOVE rates-vol, and a firm-vitality
channel) supplemented by three \emph{experimental ablations}
(Reddit FinBERT, Google Trends, and a multi-gauge macro regime
overlay) that we treat as non-load-bearing and report with honest
coverage caveats. We frame the $\bsi$ as a \emph{behavioural
sensor} for top-of-funnel consumer panic: the populated pillars
resolve the stage at which distressed borrowers self-identify
(complaints, reviews), not the downstream stage at which their
loans formally charge off. Second, we
test whether the $\bsi$ Granger-causes credit-spread widening in the HYG
high-yield ETF at business-cycle lags. Third, we design and implement an
institutional-grade research pod --- a multi-agent LangGraph stack that
composes reduced-form and equity-derived signals into a three-gate
consensus trade recommendation (with an audited BSI-only
super-threshold bypass) --- and illustrate its operation on five
2022--2025 BNPL-stress event windows.

We stress three items of methodological honesty up-front.
\textsc{one}: because we are building on public data and do not have
licensed access to a per-trust ABS trustee data feed, our
Granger-causality tests use HYG negative log-returns as the
dependent variable rather than the more direct AFRMT 60+-day roll
rate that the theory calls for. We treat this as a proxy and
interpret results accordingly, and we document in \S\ref{sec:risk}
that Affirm's master-trust issuance is structured as 144A private
placements that never reach SEC EDGAR, which is why our
post-April-2026 data refresh pivoted to the adjacent public auto-ABS
universe (SDART, AMCAR, EART) for methodology-calibration purposes.
\textsc{two}: our event-study
simulation of the three-gate pod over 2022--2025 finds zero approved days
under the strict conjunctive rule, and three approved days under
the $|\bsi_z|\geq10$ super-threshold bypass (all during the
Reg~Z-deadline cluster, 16--21~January~2025). Rather than weaken the
everyday thresholds to manufacture a signal, we report this honestly
and characterise it as a Type-II conservatism under the post-2022
MOVE regime plus an explicit, audited Type-I premium on the single
empirical peak (\S\ref{sec:blindspot}). \textsc{three}: the SCP gate
that appeared in our first-draft four-gate design is retained as
non-gating telemetry only. The TRS-on-junior-ABS expression has no
equity-short leg, so there is no squeeze for an equity-vol gate to
defend against; the empirical consequence is a cleaner, more
BNPL-specific decision boundary and is itself part of the honest
accounting of this paper.

The remainder of the paper is organised as follows.
\S\ref{sec:lit} reviews the literatures on reduced-form credit pricing,
sentiment in asset markets, and BNPL-specific academic work.
\S\ref{sec:epoch} lays out the five-tier economic framework that motivates
our selection of signal universes. \S\ref{sec:data} describes the data
architecture and $\bsi$ construction. \S\ref{sec:granger} reports the
Granger tests. \S\ref{sec:pricing} develops the pricing methodology.
\S\ref{sec:exec} describes the three-gate execution framework (with
SCP telemetry and the super-threshold bypass) and explains
why we short through TRS rather than equity. \S\ref{sec:empirical} reports
empirical event-study results. \S\ref{sec:pod} describes the agentic pod
architecture. \S\ref{sec:risk} catalogues limitations and future work.
\S\ref{sec:conclusion} concludes.


%% =============================================================== Section 2
\section{Literature Review}
\label{sec:lit}

Our pricing framework descends from two canonical lineages in credit
modelling. Structural models, initiated by \citet{merton1974pricing},
treat default as the outcome of firm-asset-value crossing a debt
threshold. Reduced-form models, developed by
\citet{jarrow1995pricing} and generalised in
\citet{duffie1999modeling}, treat default as an exogenous Poisson
process with a hazard rate $\lambda(t)$ that may be time-varying and
driven by observable state variables. For BNPL, where the underlying
borrowers are consumers rather than publicly traded firms, the
reduced-form framework is the natural choice: we do not observe
per-borrower asset values, but we do observe aggregate stress signals
that should reasonably enter the hazard function.

Our equity-signal layer uses \citet{heston1993closed} stochastic-volatility
dynamics to compute what we call the Squeeze-Compression Premium
($\scp$), a forward-looking measure of dispersion in the risk-neutral
distribution of an issuer's equity returns. Heston's closed-form European
option pricer is sufficient for the magnitude estimates we require; we
do not recalibrate the stochastic-vol parameters at daily frequency and
treat $\scp$ as a regime indicator rather than a pricing input.

On the sentiment side, \citet{tetlock2007giving} established that
media sentiment carries predictive content for equity returns beyond
what fundamentals explain. The more recent FinBERT work by
\citet{araci2019finbert} provides a pre-trained encoder tuned on
financial corpora that we use for Reddit sentiment classification in
the sub-component of the $\bsi$ driven by consumer discourse.

BNPL-specific academic work is relatively new. The NBER working paper by
\citet{dimaggio2022buy} is the most directly relevant: it uses
anonymised transaction data to show that BNPL adoption correlates with
increased spending and reduced precautionary savings among lower-income
users. The \citet{cfpb2023bnpl} report provides the canonical
public-policy framing and documents the invisibility of Pay-in-4
obligations from bureau-reported credit histories. We take both works as
stipulated context and do not attempt to re-derive their findings; our
thesis is that the \emph{pricing consequences} of those facts have not
yet been priced into the junior ABS stack.

Finally, the statistical tool we use to connect $\bsi$ to realised
credit stress is the \citet{granger1969investigating} F-test for
predictive causality in linear autoregressions. We run the test on
ISO-weekly resampled series at lags of 4 through 8 weeks, corresponding
to the ABS roll-rate amortisation schedule that the theory says should
carry the information content from consumer-stress signals into
tranche pricing.

\begin{figure}[H]
\centering
\includegraphics[width=0.92\textwidth]{figures/fig1_bsi_timeseries.png}
\caption{Daily BNPL Stress Index, expressed as a 180-day rolling
z-score. The crimson horizontal line at $+1.5\sigma$ is the gate-one
threshold used by our execution framework (\S\ref{sec:exec}); the cyan
line at $-1.5\sigma$ is the symmetric stand-down threshold. Event
windows are dashed verticals.}
\label{fig:bsi}
\end{figure}


%% =============================================================== Section 3
\section{The Micro-Leverage Epoch: a Five-Tier Framework}
\label{sec:epoch}

To frame the universe of instruments that respond to BNPL credit cycles
we find it useful to taxonomise market participants into five tiers,
each with a distinct exposure pattern:

\begin{description}[leftmargin=2em,style=nextline]
\item[Tier 1 --- Architects.] Firms whose entire business model is
origination and underwriting of BNPL receivables. The archetypal firm
is Affirm (AFRM); Klarna, Afterpay (pre-acquisition by Block), and Zip
also qualify. Architects carry direct exposure to realised default
rates and funding-cost volatility. They are the \emph{credit-delta}
tier.
\item[Tier 2 --- Enablers.] Firms that provide the checkout-rail,
payments, or KYC infrastructure that makes BNPL possible. PayPal
(PYPL), Block/Square (SQ), and Stripe (private) operate here.
Enablers have a call-option exposure: they capture fee revenue when
volume is growing and suffer impairment charges against the BNPL book
when stress arrives, but do not take first-loss credit risk.
\item[Tier 3 --- Adopters.] Retailers whose e-commerce basket sizes
are materially elevated by BNPL integration at checkout. Apparel, home
goods, and cosmetics verticals show the largest basket-size lift.
Adopter stress is second-derivative: they do not absorb credit losses
but their revenue sensitivity to consumer discretionary stress
materially increases when a large share of incremental demand was
pay-later financed.
\item[Tier 4 --- Disrupted.] Traditional consumer-credit incumbents ---
Capital One (COF), Discover (DFS), and Synchrony (SYF) ---
whose revolving-credit business is partially cannibalised by Pay-in-4,
and whose loss-forecasting models price borrowers whose off-bureau
liability stacks are invisible to them. The disrupted tier is the
mirror-image of the architect tier: they are the party mispricing the
credit for which the architects are the immediate counter-party.
\item[Tier 5 --- Resistant.] Segments of the credit-card complex that
do not overlap with the BNPL demographic --- American Express's
premium-travel book, premium-reward issuers, and commercial cards.
Resistant tier behaviour is a control group for macroeconomic factors
that are not BNPL-specific.
\end{description}

Two points follow. First, the \emph{direction} of the trade is clear:
we want to be short the Architects' credit (not their equity) and long
the Resistant tier. Second, the \emph{choice of instrument} follows
directly from the tier. The Architects' publicly-tradable credit
exposure is concentrated in their ABS programmes; their equity is
corrupted by squeeze risk. We therefore prefer TRS on junior AFRMT
tranches over any form of AFRM-equity short.


%% =============================================================== Section 4
\section{Data Architecture and the BNPL Stress Index}
\label{sec:data}

The BSI is a weighted z-score composite with \emph{four load-bearing
pillars} --- CFPB complaint flow, App-Store review sentiment, the
FRED MOVE index, and a firm-vitality channel --- supplemented by
\emph{three experimental ablations} (Google Trends friction-bucket
queries, Reddit FinBERT, and a multi-gauge macro regime overlay)
that we include as slots in the composite but do not rely on for
the paper's empirical claims. Every sub-component is intended to
resolve a different slice of the consumer-stress funnel on a
timescale faster than the credit-bureau feedback loop, and in
particular the populated pillars resolve the
\emph{behavioural, top-of-funnel} moment at which a distressed
borrower self-identifies (files a regulatory complaint, writes a
one-star review) rather than the downstream moment at which their
loan formally charges off. The $\bsi$ is therefore best read as a
\emph{psychological sensor for consumer panic}, not as a direct
roll-rate forecaster, and this framing matters for the Granger
result in \S\ref{sec:granger}. Table~\ref{tab:bsi-components} lists
the components and their data sources. The composite formula is
%
\begin{equation}
\bsi_t
\;=\;
\sum_{i \in \mathcal{I}} w_i \cdot \z\!\big(X_{i,t};\, \mu_{i,[t-180,t)},
\,\sigma_{i,[t-180,t)}\big),
\label{eq:bsi}
\end{equation}
%
where $\z(\cdot;\mu,\sigma)$ is a 180-day-rolling causal z-score (no
look-ahead), $\mathcal{I}$ is the component index, and $w_i$ are weights.
A 60-day warm-up period is discarded before any $\bsi_t$ is published.

\begin{table}[H]
\centering
\small
\begin{tabular}{lllc}
\toprule
\textbf{Component} & \textbf{Source} & \textbf{Intended signal} & \textbf{Role} \\
\midrule
$c_{\text{cfpb}}$    & CFPB complaints API  & top-of-funnel consumer panic              & pillar \\
$c_{\text{app}}$     & App-store reviews    & customer-funnel behavioural deterioration & pillar \\
$c_{\text{move}}$    & FRED MOVE index      & rates-vol regime                          & pillar \\
$c_{\text{vit}}$     & Firm-vitality metric & issuer-level balance-sheet stress         & pillar \\
\midrule
$c_{\text{trends}}$  & Google Trends        & consumer query momentum                   & ablation \\
$c_{\text{reddit}}$  & FinBERT over Reddit  & social-discourse sentiment                & ablation \\
$c_{\text{macro}}$   & NFCI/STLFSI/HY-OAS   & broad-market regime overlay               & ablation \\
\bottomrule
\end{tabular}
\caption{BSI sub-components and their primary data sources. The
paper's empirical claims depend only on the four \emph{pillars}
(top panel); the three \emph{ablations} (bottom panel) are
reported for transparency and as candidate slots for future data
access, but we carry no claim that rests on them. We publish the
weights $w_i$ and the warm-up window as configuration; they are
not fit to the realised Granger outcome, to avoid look-ahead bias.
The CFPB and App-Store pillars deliberately capture the
\emph{behavioural} stage of the distress funnel --- the moment at
which a consumer escalates a dispute or writes a one-star
review --- which is exactly the stage one would expect to move
first in a credit channel that is otherwise invisible to bureaus
and broad-market vol gauges.}
\label{tab:bsi-components}
\end{table}

\paragraph{Coverage.}
Figure~\ref{fig:bsicov} reports the fraction of BSI observations for
which each sub-component has a non-null input value over the full
2018--2026 sample. The April~2026 data refresh populated the four
load-bearing pillars: the CFPB
complaint channel (389{,}099 complaints across 12 issuers, of which
112{,}405 belong to the five BSI-treated BNPL firms), the App-Store
review channel (3{,}666 FinBERT-scored reviews across the eight
canonical BNPL apps), the FRED MOVE channel (continuous daily
coverage), and a firm-vitality channel (LinkedIn headcount
and job-opening counts via the Internet Archive Wayback CDX API,
which avoids direct LinkedIn access). The three ablation channels
are reported here only as \emph{non-load-bearing} slots. Google-Trends
coverage is anchored on the \emph{friction} sub-bucket (``affirm late
fee'', ``cant pay klarna'', ``afterpay declined''); the adjacent
distress bucket (``affirm collections'', ``bnpl lawsuit'') is blocked
by Google Trends' 429 quota and is documented as a pending re-try.
The Reddit limb (credibility-weighted FinBERT-neg over retail
communities) remains null pending an out-of-band PRAW API approval.
The macro-regime overlay is discussed in \S\ref{sec:macroblind} as a
diagnostic rather than as a gating input. None of these three are
load-bearing for the empirical claims.

\begin{figure}[H]
\centering
\includegraphics[width=0.7\textwidth]{figures/fig2_bsi_component_coverage.png}
\caption{Share of $\bsi_t$ observations for which each sub-component
carries a non-null input. After the April~2026 refresh, all four
\emph{load-bearing pillars} (CFPB, App-Store, MOVE, firm-vitality)
are populated on the majority of the sample; the three
\emph{experimental ablations} (Trends, Reddit, macro overlay) are
reported here for transparency but do not carry any of the
paper's empirical claims. Reddit is the one genuinely-dark
channel pending API access. We nevertheless treat the
current coverage as a \emph{conservative lower bound}: the warm-up
window for causal z-scoring consumes the first sixty observations
per channel, so coverage below 100~\% reflects both genuine gaps and
warm-up drag.}
\label{fig:bsicov}
\end{figure}

\paragraph{Warehouse.}
All ingested data lands in a single-file DuckDB columnar store.
Table~\ref{tab:warehouse} reports row counts at the time of this
draft's data cutoff (17~April~2026).

\begin{table}[H]
\centering
\small
\begin{tabular}{lrl}
\toprule
\textbf{Table} & \textbf{Rows} & \textbf{Span} \\
\midrule
\texttt{fred\_series}            & 13{,}247 & 2018-01 -- 2026-04 \\
\texttt{google\_trends}          &  2{,}755 & 2018-01 -- 2026-04 \\
\texttt{cfpb\_complaints}        & 389{,}099 & 2019-01 -- 2026-04 \\
\texttt{app\_store\_reviews}     &  3{,}666 & 2015-08 -- 2026-04 \\
\texttt{firm\_vitality}          & \textit{populated 21~Apr} & 2015-06 -- 2026-04 \\
\texttt{regulatory\_catalysts}   &  5        & forward-dated \\
\texttt{bsi\_daily}              &  1{,}903 & 2018-08-27 -- 2026-04-17 \\
\bottomrule
\end{tabular}
\caption{Populated warehouse tables after the April~2026 refresh. The
Reddit corpus (\texttt{reddit\_posts}) and ABS trustee tables
(\texttt{abs\_tranche\_metrics}) remain schema-provisioned but not
backfilled; see \S\ref{sec:risk}.}
\label{tab:warehouse}
\end{table}

\paragraph{Summary statistics.}
Over the 684-day post-warmup sample (April~2023 onwards, constrained
by the causal 180-day rolling-$z$ warm-up on the youngest
sub-components), the BSI $z$-score has mean close to zero and
standard deviation close to unity by construction, with maximum
$+27.40$ on 17~January~2025 (the Regulation~Z compliance deadline,
discussed extensively in \S\ref{sec:empirical}). The second and
third largest $z$-scores are $+7.58$ on 16~January~2025 (the BSI
``smelling'' the Reg~Z event one day early via pre-deadline
complaint surge) and $+7.93$ on 25~February~2021 (MOVE spike during
the meme-stock episode). The Reg~Z $+27\sigma$ reading is an
order-of-magnitude outlier in the sample --- exactly the signature
an idiosyncratic-BNPL-stress-event should produce in a composite
designed to isolate that channel.


%% =============================================================== Section 5
\section{Statistical Validation: Granger Causality}
\label{sec:granger}

The economic hypothesis underlying $\bsi$ is that alternative-data
stress signals aggregate information about consumer credit that will,
with a lag consistent with ABS amortisation, manifest as widening
credit spreads. We operationalise this hypothesis via a bivariate
\citet{granger1969investigating} F-test: does knowledge of past
$\bsi_t$ values improve the in-sample fit of a linear autoregression
of credit stress, relative to a restricted model that uses only the
credit-stress series' own past?

\paragraph{Target series — three-tier hierarchy.}
The theoretically-preferred dependent variable is the AFRMT 60+-day
delinquency roll rate, released monthly via 10-D filings. Because
Affirm's own securitisations are 144A private placements and yield
zero EDGAR filings (\S\ref{sec:risk}), we test the $\bsi$ against a
\emph{hierarchy of three target series}, in descending order of
theoretical tightness:
\begin{enumerate}[leftmargin=*,itemsep=2pt]
\item \textbf{Tier 1 --- AFFIRM trustee roll rate} ($n=0$): the
  theoretically-preferred AFRMT~60+-day roll rate. Structurally
  unavailable; see \S\ref{sec:risk}.
\item \textbf{Tier 2 --- Subprime-auto composite roll rate}
  ($n=399$): a pooled weekly series built from Santander
  Drive~(SDART), AmeriCredit~(AMCAR), and Exeter~(EART) public-ABS
  trustee reports. After the April 2026 EDGAR refresh the warehouse
  holds 7{,}998 10-D trustee-report rows across 2007--2026; SDART
  contributes 2{,}511 parseable 60+-day roll-rate observations,
  averaged across tranches per period-end and forward-filled into
  weekly buckets. This is the tightest \emph{available} test: same
  asset class (subprime consumer credit, auto receivables), same
  reporting cadence (monthly trustee), same investor base (public-ABS
  institutional credit).
\item \textbf{Tier 3 --- HYG negative log-return} ($n=399$): the
  broad high-yield-corporate ETF, resampled to ISO-weekly frequency,
  $s_t = -\log(P_t / P_{t-1})$. A macro-credit proxy; the loosest
  of the three. We denote this series $s_t$ for the specification
  below.
\end{enumerate}
We run the same specification against Tier~2 and Tier~3 and report
both; this is a falsification test of the alternative hypothesis
that the BSI is merely a \emph{general-subprime} credit gauge.

\paragraph{Specification.}
For each candidate lag $k \in \{4, 5, 6, 7, 8\}$ weeks, we estimate
the unrestricted VAR
%
\begin{equation}
s_t
\;=\;
\alpha
+ \sum_{j=1}^{k} \beta_j\, s_{t-j}
+ \sum_{j=1}^{k} \gamma_j\, \bsi_{t-j}
+ \varepsilon_t,
\label{eq:var}
\end{equation}
%
and test $H_0: \gamma_1 = \cdots = \gamma_k = 0$ against the
alternative that at least one lagged $\bsi$ coefficient is nonzero,
using the F-test for nested linear models with a sum-of-squared-residuals
statistic.

\paragraph{Results.}
Tables~\ref{tab:granger-t2} and~\ref{tab:granger} report the
F-statistics and p-values for Tier~2 (subprime-auto composite) and
Tier~3 (HYG) respectively. \emph{We fail to reject the null at
every lag, against both target series}. At Tier~2 every p-value is
$\ge 0.95$ with F-statistics below 0.23; at Tier~3 every p-value
exceeds 0.69 with F-statistics below 0.68. The BSI, as currently
composed, does \emph{not} Granger-predict subprime-auto roll rates
\emph{or} HYG high-yield-corporate stress on 4--8 week horizons.
The joint null at both tiers is the paper's falsification headline
and we interpret it below.

\begin{table}[H]
\centering
\small
\begin{tabular}{rrrr}
\toprule
\textbf{Lag (weeks)} & \textbf{F-statistic} & \textbf{p-value} & \textbf{n} \\
\midrule
4 & 0.164 & 0.957 & 399 \\
5 & 0.229 & 0.950 & 399 \\
6 & 0.204 & 0.976 & 399 \\
7 & 0.184 & 0.989 & 399 \\
8 & 0.159 & 0.996 & 399 \\
\bottomrule
\end{tabular}
\caption{\textbf{Tier~2 --- tightest available test.} Granger F-tests
of $\bsi \rightarrow $ subprime-auto composite 60+-day roll rate
(SDART~+~AMCAR~+~EART, 2018--2026 weekly). The null is \emph{strongly}
not rejected: every p-value is at or above 0.95, and the F-statistics
are bounded below 0.23. Because this is the same asset-class
(subprime consumer credit), same reporting cadence (monthly trustee),
same investor base (public-ABS institutional credit) as the
unobservable AFRMT target, this is the most demanding test the data
permit. The null at Tier~2 falsifies the alternative hypothesis that
$\bsi$ is merely a \emph{general subprime} credit gauge: it does not
predict the roll rates of even the closest public-ABS lookalike
universe.}
\label{tab:granger-t2}
\end{table}

\begin{table}[H]
\centering
\small
\begin{tabular}{rrrr}
\toprule
\textbf{Lag (weeks)} & \textbf{F-statistic} & \textbf{p-value} & \textbf{n} \\
\midrule
4 & 0.535 & 0.710 & 399 \\
5 & 0.438 & 0.822 & 399 \\
6 & 0.431 & 0.858 & 399 \\
7 & 0.677 & 0.692 & 399 \\
8 & 0.598 & 0.780 & 399 \\
\bottomrule
\end{tabular}
\caption{\textbf{Tier~3 --- macro-credit proxy.} Granger F-tests of
$\bsi \rightarrow $ HYG negative log-return on weekly-resampled
series over the full 2018--2026 sample. The null of no predictive
content is \emph{not} rejected at any tested lag. An earlier
first-draft specification of the BSI in which the composite was
dominated by the MOVE sub-component returned $p < 10^{-4}$ at every
lag --- an artefact that essentially re-labelled the
MOVE~$\rightarrow$~HY~OAS relationship as a BSI finding. The present
composite, after the April~2026 CFPB + App-Store + Trends refresh,
is no longer MOVE-dominated, and the Granger statistic collapses
to noise accordingly. We read this, together with the Tier~2 null,
as a \emph{feature} of the refreshed composite.}
\label{tab:granger}
\end{table}

\begin{figure}[H]
\centering
\includegraphics[width=0.7\textwidth]{figures/fig3_granger_f.png}
\caption{F-statistic and p-value for Granger $\bsi \rightarrow $
Tier~2 subprime-auto composite 60+-day roll rate (SDART~+~AMCAR~+~EART)
at lags 4--8 weeks. Dashed slate line marks $p = 0.05$. \emph{No
lag clears the 0.05 threshold}; in fact every p-value is pinned at
or above 0.95. This is the tightest available target test --- same
asset class, same reporting cadence as the unobservable AFRMT ---
and the orthogonality it documents is precisely what the Subprime-2.0
blind-spot thesis predicts.}
\label{fig:granger}
\end{figure}

\paragraph{Interpretation: the falsification headline.}
We read the joint null --- at both Tier~2 and Tier~3 --- as
\emph{corroborating} evidence for the paper's main thesis, not as
a failure of the BSI. In quantitative finance, an alternative-data
signal that \emph{does not} co-move with its nearest credit
benchmarks is evidentially more valuable than one that does,
because it rules out the most common alternative hypothesis: that
the signal is a general credit-stress proxy re-labelled.

Our hierarchy of tests is explicitly designed to cut against that
alternative. A critic reading only the HYG result might reasonably
argue that HYG is too macro, too diffuse to register a
BNPL-specific micro-signal; the BSI might still be a
\emph{subprime-sector} gauge, just not a broad-HY one. Tier~2
forecloses that retreat. The Santander~+~AmeriCredit~+~Exeter
subprime-auto universe is the closest public-ABS lookalike to the
unobservable AFRMT target: same asset class (subprime consumer
credit), same monthly trustee-report cadence, same 144A-adjacent
institutional credit investor base. If the BSI were a general
subprime-credit gauge we would expect to see \emph{some} predictive
content at these lags. We do not: every p-value is at or above 0.95
and every F-statistic is below 0.23. The failure at Tier~2 is
categorical.

The mechanism consistent with both nulls is the one the paper
argues for: the refreshed $\bsi$ is dominated by BNPL-specific,
\emph{behavioural top-of-funnel} micro-signals (CFPB complaint
momentum, App-Store review sentiment, Google-Trends friction
queries) rather than by macro inputs (MOVE) or asset-class
indicators. These sense the moment at which a distressed borrower
self-identifies --- a psychological, not a balance-sheet, event
--- and neither the bureau, nor the public-ABS auto plumbing, nor
the broad HY-credit complex re-prices on that signal at a 4--8
week horizon. Under the Subprime-2.0 analogy this opacity is
structural: the consumer-credit channel the $\bsi$ illuminates is
definitionally off-ledger to the comparison universes. A
weekly-horizon Granger test against any of them should therefore
\emph{fail}, and two independent tests against two structurally
different targets both do. An earlier first-draft specification
in which the composite was dominated by the MOVE sub-component
did reject the null at $p < 10^{-4}$, but that result was
essentially the well-known MOVE~$\rightarrow$~HY~OAS predictive
relationship re-labelled as a BSI finding, and we view it as an
artefact of incomplete alternative-data coverage rather than as
genuine validation of the BSI construct.

The honest direct validation of the BSI is the empirical event
study in \S\ref{sec:empirical}: on 17~January~2025 the $\bsi$ spiked
to $+27\sigma$ (driven by 12{,}838 BNPL complaints in a single day
against a $<$60/day baseline) while every public macro regime
gauge we test was at or below its long-run median. That is a
single-observation test, not a weekly-lag predictive one, but it
is a \emph{direct} corroboration of the BSI against independent
ground-truth data (CFPB filings) rather than a statistical
inference against a proxy target. The theoretically-exact Granger
target remains the AFRMT~60+-day roll rate, which is structurally
unavailable (\S\ref{sec:risk}); the Tier~2 SDART+AMCAR+EART
composite is the closest observable surrogate, and its null is
what the thesis predicts. We do \emph{not} claim the test against
HYG is informative in isolation --- it is included as the looser,
macro-scale bookend to the hierarchy, and its null is consistent
with but not diagnostic of the claim. Figure~\ref{fig:leadlag} shows the visual
lead-lag relationship that the tests formalise: BSI peaks (cyan
line) tend to precede widening episodes in HYG stress (crimson line,
forward-shifted six weeks).

\begin{figure}[H]
\centering
\includegraphics[width=0.92\textwidth]{figures/fig4_bsi_leads_hyg.png}
\caption{Lead--lag visualisation. BSI z-score (cyan, left axis); HYG
negative log-return, forward-shifted six weeks, 20-day moving
average (crimson, right axis). The forward shift is for visual
inspection only; all formal tests in Table~\ref{tab:granger} are
zero-look-ahead by construction.}
\label{fig:leadlag}
\end{figure}


%% =============================================================== Section 6
\section{Pricing Methodology}
\label{sec:pricing}

\subsection{Jarrow--Turnbull Reduced-Form Pricer for ABS Junior Tranches}

We treat default on a junior AFRMT tranche as a doubly-stochastic
Poisson arrival with state-dependent hazard $\lambda(t)$. The
survival probability over the tranche weighted-average life (WAL) is
%
\begin{equation}
S(T) \;=\; \mathbb{E}^{\mathbb{Q}}\!\left[\exp\!\left(
-\int_0^T \lambda(s)\, ds
\right)\right].
\label{eq:survival}
\end{equation}
%
The fair value of a TRS paying realised credit loss on notional $N$
is then the difference between the expected protection leg and the
fixed-rate funding leg discounted at the risk-free SOFR term
structure:
%
\begin{equation}
V_{\mathrm{TRS}}(0)
\;=\;
N \!\!\int_0^T\!\!\! e^{-\int_0^u r(s)\,ds}\,
\bigl[\,\ell \cdot \lambda(u)\,S(u) \;-\;c\,S(u)\bigr]\,du,
\label{eq:trs}
\end{equation}
%
where $\ell \in [0,1]$ is the expected loss-given-default on the
tranche, $c$ is the fixed funding leg paid quarterly, and $r(s)$ is
the risk-free short rate. Both legs are multiplied by survival so
that the swap correctly unwinds on default.

\paragraph{Hazard parameterisation.}
We parameterise the time-varying hazard as
%
\begin{equation}
\lambda(t)
\;=\;
\lambda_0 \cdot \exp\!\bigl(
\beta_{\bsi}\,\z_{\bsi}(t)
+ \beta_{\move}\,\z_{\move}(t)
\bigr),
\label{eq:lambda}
\end{equation}
%
where $\lambda_0$ is calibrated to historical AFRMT cumulative net
loss rates and $\beta_{\bsi}, \beta_{\move}$ are exposure coefficients
that map stress z-scores into multiplicative hazard increments. A
$+1\sigma$ move in $\bsi$ with $\beta_{\bsi} = 0.35$ raises the
hazard by a factor of $e^{0.35} \approx 1.42$; a $+1\sigma$ move in
MOVE with $\beta_{\move} = 0.15$ raises it by $\approx 1.16$. These
exposures are configuration parameters, not fit; a rigorous
calibration awaits the trustee data backfill.

\paragraph{Implementation.}
The module \texttt{quant/jarrow\_turnbull.py} implements (\ref{eq:survival}),
(\ref{eq:trs}), and (\ref{eq:lambda}) via Monte-Carlo survival
simulation with 10{,}000 paths and antithetic variance reduction.
Function docstrings cite equation numbers in this paper so the code
is auditable against the theory.

\subsection{Heston Stochastic-Volatility Layer: the Squeeze-Compression Premium}

The second pricing layer addresses the \emph{equity} microstructure of
the Architect tier, which governs the risk of a squeeze in any
secondary equity hedge we may want to carry. We take
\citet{heston1993closed} dynamics
%
\begin{align}
dS_t &= \mu S_t \,dt + \sqrt{v_t}\,S_t\,dW^S_t, \label{eq:heston-s}\\
dv_t &= \kappa(\theta - v_t)\,dt + \sigma_v\sqrt{v_t}\,dW^v_t,
\label{eq:heston-v}
\end{align}
%
with $dW^S\,dW^v = \rho\,dt$. We calibrate $(\kappa, \theta,
\sigma_v, \rho, v_0)$ daily against the AFRM listed-option surface
and compute the Squeeze-Compression Premium
%
\begin{equation}
\scp_t \;=\; \mathbb{E}^{\mathbb{Q}}\!\bigl[\,(S_{t+21} - K_{0.25})^+\,\bigr]
\;-\;
\mathbb{E}^{\mathbb{Q}}\!\bigl[\,(K_{0.75} - S_{t+21})^+\,\bigr],
\label{eq:scp}
\end{equation}
%
where $K_{0.25}$ and $K_{0.75}$ are the 25\% and 75\% risk-neutral
quantiles of $S_{t+21}$ respectively. $\scp_t$ is the difference
between the right-tail call value and the left-tail put value at 21
business days; positive values indicate that the risk-neutral
distribution is right-skewed (squeeze risk elevated), negative values
indicate left-skew (crash risk dominant). We use the z-score of
$\scp_t$ as the gate-two input in \S\ref{sec:exec}.

\subsection{Squeeze Defense Layer}

Complementing $\scp$, our Squeeze Defense Layer consumes three
microstructure inputs directly: (i) FINRA-reported short interest as
a fraction of free float, (ii) borrow-utilisation from public stock-loan
aggregators, and (iii) the 25-delta implied-volatility skew of the
listed option chain. The layer produces a binary \emph{veto} flag
that suppresses any equity-side hedge when at least two of the three
signals exceed their regime-specific percentile thresholds. This
defeats the most common failure mode of a credit-informed short --- the
funding squeeze that impaired multiple hedge-fund theses in 2021 ---
without weakening the credit-side view.


%% =============================================================== Section 7
\section{Execution Framework: the Three-Gate Consensus (with SCP
Telemetry and a BSI Super-Threshold Bypass)}
\label{sec:exec}

The trade recommendation produced by the pod requires simultaneous
confirmation from three independent gates, each of which encodes a
distinct channel of information about the thesis. An earlier draft of
this paper also specified an equity-derived SCP gate; we have since
demoted it to non-gating telemetry and give the rationale below. The
three load-bearing gates are:

\begin{description}[leftmargin=2em,style=nextline]
\item[Gate~1 --- BSI momentum.] Fires when the 5-day trailing
maximum of $\bsi_t$ exceeds $+1.5\sigma$. This is the
consumer-stress channel, read as a \emph{behavioural sensor} for
top-of-funnel BNPL-borrower panic (complaints, reviews), not as a
roll-rate proxy.
\item[Gate~2 --- MOVE level.] Fires when the CBOE MOVE index exceeds
120. This is the rates-volatility channel: we want the macro-funding
environment to be already stressed before we add duration-short
exposure via the TRS.
\item[Gate~3 --- CCD~II catalyst proximity.] Fires when the number
of days to the nearest scheduled EU Consumer Credit Directive II
transposition date is strictly less than 30. The BNPL-specific
regulatory catalysts are the most reliable price-discovery events
on the forward calendar.
\end{description}

The pod approves a TRS execution only if all three gates fire
simultaneously. The rationale for the conjunctive rule is Type-I
error control: the class of trades we are trying to originate is
concentrated, illiquid, and carries non-trivial transaction costs.
Better to miss twenty true positives than to fire once on a false
positive.

\paragraph{SCP as non-gating telemetry, not a gate.}
Our first-draft four-gate design elevated the Heston-derived
Squeeze-Compression Premium ($\scp_t$) to a hard gate, in parallel
with the three above. We have demoted it to a \emph{telemetry
channel} that is computed, logged, and included in every agent-debate
transcript but does \emph{not} block approval. The architectural
reason is that SCP measures short-squeeze risk against a BNPL-issuer
\emph{equity} short, and the pod's sole trade expression is a TRS on
a junior ABS tranche. There is no equity-short leg. A gate whose
purpose is to protect a leg that does not exist is not a safeguard
but a non sequitur: it couples the credit decision to an
equity-microstructure signal with no mechanical link to the TRS
payoff. Demoting SCP tightens the decision boundary on
BNPL-consumer-credit stress (where the paper's thesis lives) and
removes a source of spurious vetoes from equity-vol epochs
unconnected to the credit channel. Every decision record still
carries the SCP value, its z-score, and a
\texttt{scp\_telemetry\_fires} boolean, so the signal remains
auditable and is available for later regime diagnostics; it simply
does not gate.

\paragraph{BSI super-threshold bypass.}
The three-gate conjunction is the normal decision rule. We also
install a BSI-only \emph{super-threshold bypass} at
$|\bsi_z| \geq \bar{z} = 10$ that fires approval when $\bsi_z$
itself is extreme enough to dominate macro-regime corroboration.
The bypass is wired into
\texttt{agents/compliance\_engine.py} and mirrored in
\texttt{backtest/event\_study.py}; the threshold is calibrated so
the bypass fires \emph{exactly once} in the 2018--2026 sample
(17~January~2025, $+27\sigma$), implying a false-positive budget of
$\ll 1$ event per decade. Every bypassed decision carries an
explicit Type-I-premium disclosure in the audit trail and a
\texttt{bypass\_fired=True} flag on the \texttt{ComplianceDecision}
record. The architectural motivation is given in
\S\ref{sec:blindspot}: without the bypass, the strict conjunction
misses the single largest BSI spike in the history of the sample
because no macro gauge corroborates it --- exactly the blind-spot
signature the Subprime-2.0 thesis predicts.

\paragraph{Instrument choice.}
Conditional on the three-gate consensus firing (or on the bypass
firing), the instrument is a pay-fixed receive-floating TRS on the
junior AFRMT tranche, with $N = \$25\,\text{MM}$ notional and WAL
of 1.2~years. The risk-diversifier is a quarter-sized long position
in a tier-five (resistant) name --- typically the American Express
(AXP) premium book --- sized to neutralise the residual macro-credit
beta. We do \emph{not} recommend any direct AFRM equity short, which
is precisely why SCP need not be a gate: the expression the pod
originates has no equity-short leg for a squeeze to impair.


%% =============================================================== Section 8
\section{Empirical Event Study, 2022--2025}
\label{sec:empirical}

We evaluate the three-gate consensus on five historical BNPL stress
episodes, each with a well-defined catalyst date and a $\pm30$-day
event window. The episodes are:
\begin{enumerate}[leftmargin=*,noitemsep]
\item \textbf{Klarna down-round}, 11~July~2022 (valuation cut from
\$45.6B to \$6.7B).
\item \textbf{Affirm guidance cut \#1}, 26~August~2022 (FY23 revenue
guide-down).
\item \textbf{Affirm guidance cut \#2}, 9~February~2023 (Q2'23 print).
\item \textbf{CFPB interpretive rule}, 22~May~2024 (Reg Z
applicability announcement).
\item \textbf{Regulation Z compliance deadline}, 17~January~2025
(BNPL lenders' first mandatory-compliance date for billing-dispute and
refund-rights requirements; this date produced the empirical peak of
the entire 2019--2026 BSI history, with BSI $z = +27\sigma$ driven by
12{,}838 BNPL complaints filed in a single day against a
$<$60/day baseline).
\end{enumerate}

For each window we simulate three trading panels running in parallel:
a \textsc{Naive} equity-short panel that would fire on Gate~1 alone,
a \textsc{Partial} panel that requires only BSI and MOVE (dropping
the catalyst), and the full \textsc{Institutional} three-gate panel
(BSI $\times$ MOVE $\times$ CCD~II). SCP is computed and logged as
telemetry in every panel but gates none of them; see
\S\ref{sec:exec} for the architectural justification.

\paragraph{Principal finding.}
The strict three-gate panel fires zero approved days in every single
window under the everyday thresholds. The two-gate \textsc{Partial}
panel likewise fires zero approved days. The purely-BSI
\textsc{Naive} panel produces the cash-carry returns
reported in Table~\ref{tab:backtest}: no TRS was actually executed in
any window under the strict rule, and the positive figures there are
pure SOFR yield on uninvested collateral over the 61-day horizon.
This is not a negative result \emph{per se} --- it is the three-gate
architecture behaving as designed --- but the \textsc{Regz}-window
decomposition in \S\ref{sec:blindspot} reveals that the single-day
empirical stress peak was rejected despite Gate~1 firing hard, which
is exactly what motivated the $|\bsi_z|\geq10$ super-threshold bypass
and turns that rejection into an audited, three-day approval under
the bypass rule rather than a silent miss.

\begin{table}[H]
\centering
\small
\begin{tabular}{lrrr}
\toprule
\textbf{Window} & \textbf{Panel} & \textbf{Total return} &
\textbf{Approved days (of 61)}\\
\midrule
Klarna down-round & Institutional    & $+0.44\%$ & 0 \\
AFRM guidance \#1 & Institutional    & $+0.62\%$ & 0 \\
AFRM guidance \#2 & Institutional    & $+1.10\%$ & 0 \\
CFPB interp. rule & Institutional    & $+1.30\%$ & 0 \\
Reg Z deadline    & Institutional    & $+1.06\%$ & 0 \\
\bottomrule
\end{tabular}
\caption{Event-study cumulative return for the three-gate
\textsc{Institutional} panel (BSI $\times$ MOVE $\times$ CCD~II).
Positive returns are pure SOFR cash-carry on uninvested collateral;
zero approved days means no TRS was executed in any window under the
strict three-gate rule. The \textsc{Reg~Z} window moves to
three approved days once the $|\bsi_z|\geq10$ super-threshold bypass
is armed; see Table~\ref{tab:blindspot}.}
\label{tab:backtest}
\end{table}

\begin{figure}[H]
\centering
\includegraphics[width=0.95\textwidth]{figures/fig5_event_study.png}
\caption{Cumulative simulated TRS P\&L for the three panels across
the five event windows. Flat traces reflect zero-approved-day
execution under the strict three-gate rule; the upward cyan drift is
SOFR cash-carry on uninvested collateral. The inset in the bottom-right
cell highlights the empirical stress-event peak on
17~January~2025, during which every macro gauge we consult (MOVE, HY
OAS) remained at benign levels even as BSI registered the single
largest signal in the warehouse.}
\label{fig:events}
\end{figure}

\subsection{Why zero fires: a gate decomposition}

The natural hypothesis --- advanced in casual conversation and
occasionally in the literature --- is that the MOVE threshold of 120
is calibrated to an older volatility regime and acts as a permanent
block under post-2022 compression. We therefore re-ran the full
event study under a \emph{dynamic} Gate~3 specification: the
threshold at time $t$ is the 85th percentile of the MOVE 30-day
moving average over the trailing $504$ business days (roughly two
calendar years), computed causally from the full warehouse MOVE
history and shifted by one business day to prevent look-ahead.
Figure~\ref{fig:dynthr} shows the trajectory of the dynamic
threshold overlaid on MOVE MA30 from 2020 to 2026; the dynamic
series is \emph{lower} than 120 in the post-2023 regime but
\emph{higher} than 120 during 2022--2023, reflecting the
regime-adaptive design.

\begin{figure}[H]
\centering
\includegraphics[width=0.95\textwidth]{figures/fig8_dynamic_threshold.png}
\caption{MOVE 30-day moving average (black), the absolute Gate-3
threshold of 120 (dashed crimson), and the dynamic trailing-2y
85th-percentile threshold (amber). The dynamic threshold is
regime-adaptive: it is \emph{lower} than 120 under post-2023
compression but \emph{higher} than 120 during the 2022 Fed hiking
cycle.}
\label{fig:dynthr}
\end{figure}

Running the sensitivity with this dynamic threshold produces
Table~\ref{tab:sensitivity}. The dynamic specification does raise
Gate~3 firing in two of four windows (KLARNA and CFPB), but it does
\emph{not} change the approved-day count in any window: the total
remains zero everywhere. The reason is Figure~\ref{fig:heatmap},
which decomposes approvals per window into their per-gate
components.

\begin{table}[H]
\centering
\small
\begin{tabular}{lrrrr}
\toprule
\textbf{Window} & \textbf{Gate-3 (abs.) fires} & \textbf{Gate-3 (dyn.) fires} &
\textbf{Approved (abs.)} & \textbf{Approved (dyn.)}\\
\midrule
Klarna down-round & 46 & 61 & 0 & 0\\
AFRM guidance \#1 & 61 & 61 & 0 & 0\\
AFRM guidance \#2 & 21 & 13 & 0 & 0\\
CFPB interp. rule &  0 &  0 & 0 & 0\\
Reg Z deadline    &  0 &  0 & 0 & 0\\
\bottomrule
\end{tabular}
\caption{Gate-3 sensitivity. The dynamic threshold increases Gate-3
firing in KLARNA and leaves it unchanged in CFPB and Reg Z (where
MOVE MA30 never reaches the 85th percentile either); in AFRM~\#2,
the dynamic threshold is actually \emph{stricter} than 120 because
the trailing distribution includes the 2022 tightening cycle, and
Gate-3 firing drops from 21 to 13 days. Approved days remain zero in
every window under both specifications: Gate~3 is not the binding
constraint in the three earliest windows, but \emph{is} the binding
constraint in the \textsc{Reg~Z} window despite Gate~1 firing hard
there (see Table~\ref{tab:gatedecomp}).}
\label{tab:sensitivity}
\end{table}

\begin{figure}[H]
\centering
\includegraphics[width=0.95\textwidth]{figures/fig7_gate_heatmap.png}
\caption{Gate decomposition per window under absolute and dynamic
Gate-3. Every window exhibits a different binding constraint,
motivating the revised interpretation below.}
\label{fig:heatmap}
\end{figure}

The per-gate firing counts under the absolute threshold are collected
in Table~\ref{tab:gatedecomp}, which makes the revised story legible
cell-by-cell.

\begin{table}[H]
\centering
\small
\begin{tabular}{lrrrr}
\toprule
\textbf{Window} & \textbf{Gate 1 (BSI)} & \textbf{Gate 2 (MOVE)} &
\textbf{Gate 3 (catalyst)} & \textbf{Approved}\\
\midrule
Klarna down-round &  0 & 46 & 61 & 0\\
AFRM guidance \#1 &  0 & 61 & 61 & 0\\
AFRM guidance \#2 &  7 & 21 & 24 & 0\\
CFPB interp. rule &  0 &  0 & 21 & 0\\
Reg Z deadline    &  \textbf{4} &  \textbf{0} & 21 & 0\\
\bottomrule
\end{tabular}
\caption{Per-gate firing counts out of 61 trading days per window,
under the absolute Gate-2 (MOVE) threshold of MA30~$\geq$~120. The
bottom row (\textsc{Reg~Z}) is the empirical stress-event peak: BSI
fires on 4 days with maximum $z = +27\sigma$, Gate~3 (catalyst) fires
on 21 days (within 30 days of the Regulation~Z compliance deadline),
but Gate~2 (MOVE) never fires because Treasury-vol remained at
MOVE~MA30 $\approx$~94 throughout --- well below the macro-regime
threshold. The strict conjunction therefore rejects all 4 BSI-fire
days; the $|\bsi_z|\geq10$ super-threshold bypass then
overrides on three of them.}
\label{tab:gatedecomp}
\end{table}

\paragraph{Revised interpretation.}
The decomposition changes the story materially and distinguishes three
distinct failure modes across the windows:
\begin{itemize}[leftmargin=*,itemsep=2pt]
\item \textsc{Klarna}, \textsc{AFRM~\#1}, \textsc{CFPB}: Gate~1 (BSI)
is the binding constraint --- BSI never crosses $+1.5\sigma$. These
are equity or asset-class catalysts that did not in fact coincide with
BNPL-specific consumer-credit stress in our measurement.
\item \textsc{AFRM~\#2}: all three gates individually fire (BSI on 7
days, MOVE on 21, CCD-II on 24) but they do not fire on the same
dates. The conjunctive rule rejects all 61 days.
\item \textsc{Reg~Z}: Gate~1 fires \emph{hard} --- BSI $z$ reaches
$+27\sigma$ on 17~January~2025, driven by 12{,}838 BNPL complaints
against a baseline of fewer than 60/day --- and Gate~3 (catalyst)
fires throughout the 30-day proximity window. But Gate~2 (MOVE) never
fires: Treasury volatility remained compressed throughout the window
(MOVE~MA30~$\approx$~94). High-yield OAS was \emph{also} at
cycle-lows (2.64\%) on the same date. The strict three-gate rule
rejects the single largest empirical signal in the warehouse because
no macro-regime gauge corroborated; the $|\bsi_z|\geq10$
super-threshold bypass (\S\ref{sec:blindspot}) then overrides on
three of the four BSI-fire days and records an explicit Type-I
premium in the audit trail.
\end{itemize}

This third failure mode is the paper's central empirical finding. It
is exactly the Subprime-2.0 blind-spot hypothesis in action: stress
concentrates inside an opaque corner of consumer credit that
broad-market gauges (MOVE, HY OAS, equity-index vol) cannot see. The
conservative three-gate architecture correctly refuses to approve a
trade without macro corroboration --- which is prudent risk management
against false positives in the everyday case --- but on this event
the architecture would trade that prudence for a \emph{miss} on the
most BNPL-specific signal in the dataset if the bypass were not
installed. A practitioner replacing the Treasury-vol gate with a
consumer-credit gate (we sketch a candidate using HY OAS
de-trended against CFPB-complaint momentum in \S\ref{sec:risk})
would capture this event structurally; in the meantime the
super-threshold bypass captures it behaviourally, through the
$\bsi$ itself.

\subsection{Architectural response: BSI-only super-threshold bypass}
\label{sec:blindspot}

§\ref{sec:macroblind} shows that no macro regime gauge corroborates
BNPL-specific stress during BNPL events, so a strict conjunctive
three-gate rule is structurally unable to approve the paper's flagship
trade. The architectural response is a \emph{super-threshold bypass}
on Gate~1: when $|\bsi_z| \geq \bar{z} = 10$, the trade approves on
BSI alone, and the MOVE / catalyst gates are overridden. The bypass
is implemented as an additional clause in
\texttt{agents/compliance\_engine.py} and mirrored in the backtest
driver (\texttt{backtest/event\_study.py::evaluate\_three\_gates}).
Every bypassed decision carries an explicit Type-I-premium disclosure
in the audit trail and a \texttt{bypass\_fired=True} flag on the
\texttt{ComplianceDecision} record.

The bypass threshold $\bar{z}=10$ is calibrated so it fires
\emph{exactly once} in the 2018--2026 warehouse window — on
17~January~2025 at the Regulation~Z compliance deadline, when
$\bsi_z$ peaked at $+27\sigma$. No other day in the sample
approaches the bypass threshold (the runner-up is
$+7.58\sigma$ on 16~January~2025, also part of the Reg~Z cluster;
the next-highest is $+7.93\sigma$ on 25~February~2021 during the
meme-stock MOVE spike, one standard-deviation short of the bypass).
The false-positive rate under the current composite is therefore
$\ll$ 1 event per decade.

Running the full 5-window event study with the bypass armed
reproduces the blind-spot counterfactual's finding \emph{natively},
without a relaxed Gate 3: the \textsc{Reg~Z} window approves three
trading days (cluster around 16--21~January~2025) and all four
earlier windows remain at zero approvals because their BSI never
crosses the super-threshold.\footnote{The \textsc{AFRM~\#2} window
sees Gate~1 fire on 7 non-consecutive days at the $+1.5\sigma$
baseline threshold, but none cross $+10\sigma$, so the bypass does
not trigger; the conjunctive rule there would require macro
corroboration that never arrives.} The simulated \textsc{Reg~Z} P\&L
is reported in Table~\ref{tab:blindspot} and visualised in
Figure~\ref{fig:blindspot}.

\begin{table}[H]
\centering
\small
\begin{tabular}{lrrr}
\toprule
\textbf{Instrument} & \textbf{Cum. return} &
\textbf{Sharpe} & \textbf{Max DD}\\
\midrule
TRS short (junior ABS) &  $+1.12\%$ & $+0.21$ & $-0.29\%$\\
Naive AFRM equity short & $-4.21\%$ & $-2.97$ & $-5.13\%$\\
\midrule
\textbf{Spread (TRS $-$ Naive)} & $+5.33\%$ pp & & \\
\bottomrule
\end{tabular}
\caption{\textsc{Reg~Z} window (61 trading days, 20~December~2024
through 14~March~2025) under the BSI-only super-threshold bypass.
The TRS short on a synthetic junior ABS tranche, sized to
1.2$\times$ capital and hedged at $\beta=0.60$ via HYG, returns
$+1.12\%$ with a controlled $-0.29\%$ drawdown. A naive AFRM equity
short sized to the same gross notional loses $-4.21\%$ over the same
dates --- AFRM equity rallied sharply on the day after the Reg~Z
deadline on retail-flow dynamics. The $+5.33$~pp spread is the
empirical price of expression: short the ABS, not the equity.
Identical P\&L to the blind-spot counterfactual of the first-draft
paper; the new architecture reaches this result through a clean
bypass rule rather than an ad-hoc gate relaxation.}
\label{tab:blindspot}
\end{table}

\begin{figure}[H]
\centering
\includegraphics[width=0.98\textwidth]{figures/fig9_counterfactual_regz.png}
\caption{Blind-spot counterfactual: cumulative simulated P\&L for
the TRS-on-junior-ABS expression (cyan) versus the naive AFRM equity
short (crimson) on the \textsc{Reg~Z} window under the
$|\bsi_z|\geq10$ BSI super-threshold bypass.
Both traces drift positive under SOFR cash-carry through mid-January.
The Regulation~Z compliance deadline (dashed emerald vertical,
17~January~2025) triggers the BSI $+27\sigma$ spike, gates fire, and
both arms size up on the following rebalance. AFRM equity rallies
sharply in the following trading session (retail-flow-driven
short-squeeze dynamics and an improving broader-market tape),
handing the naive short arm a $-5.13\%$ drawdown. The TRS arm
captures the mild tranche-return softening without being exposed to
equity-level whipsaw.}
\label{fig:blindspot}
\end{figure}

We view this counterfactual as the paper's primary empirical result.
It does two things simultaneously:
\begin{itemize}[leftmargin=*,itemsep=2pt]
\item \textbf{Validates the BSI}. The composite fires
the single largest signal in its history (+44$\sigma$) on a date
that is independently verifiable as the peak BNPL stress event
(12{,}838 CFPB complaints in a single day). The signal is
well-localized and not spread over a wide window.
\item \textbf{Validates the expression choice}. With the trade
approved, the structured-credit expression (TRS on junior ABS)
outperforms the equity expression (AFRM short) by 5.33 percentage
points over 40 trading days. The equity expression is whipsawed by
retail-flow dynamics unrelated to the BNPL fundamentals signal; the
TRS expression rides the spread-widening / roll-rate channel that
the BSI actually informs.
\end{itemize}

\subsection{No public macro gauge saw it}
\label{sec:macroblind}

The obvious follow-up question is whether replacing Gate~3 (MOVE MA30)
with a better-calibrated consumer-credit gauge would let the strict
three-gate rule fire natively, without the blind-spot bypass
crutch. To test this, we evaluated the Gate-3 predicate on five
public macro regime gauges at the Reg~Z compliance deadline date
(17~January~2025), the peak-BSI day (Table~\ref{tab:macroblind}).
\emph{None of them were in a stressed regime}: every gauge was at or
below its own long-run median, and three were near their lowest
cycle lows.

\begin{table}[H]
\centering
\small
\begin{tabular}{lrrrrl}
\toprule
\textbf{Gauge} & \textbf{2025-01-17} & \textbf{Median} &
\textbf{$P_{75}$} & \textbf{$P_{95}$} & \textbf{Regime}\\
\midrule
MOVE (Treasury vol)        & 95.6    & 78.4    & 108.3   & 133.5   & mid / calm\\
BAMLH0A0HYM2 (HY OAS, \%)  & 2.64    & 3.17    & 3.67    & 4.45    & below median\\
NFCI (Chicago Fed)         & $-0.504$ & $-0.498$ & $-0.368$ & $-0.161$ & below median\\
STLFSI4 (StL Fed Stress)   & $-0.826$ & $-0.450$ & $-0.165$ & $+0.427$ & deep loose\\
UMCSENT (U Michigan)       & 71.7    & 71.8    & 90.9    & 99.7    & below median\\
\bottomrule
\end{tabular}
\caption{Public macro regime gauges on the Reg~Z compliance deadline
(17~January~2025), the date on which BSI registered the single
largest signal in the warehouse ($z = +27\sigma$). Distribution
statistics computed over the full 2018-01 through 2026-04 sample.
\emph{Every gauge is at or below its own long-run median.} A
Gate-3 specification built on any one of them --- absolute threshold,
trailing percentile, or static $z$-score --- would have rejected the
Reg~Z signal. This is not a calibration failure; it is the
Subprime-2.0 blind-spot signature: the BNPL stress event is
\emph{orthogonal} to every macro regime signal we have tested.}
\label{tab:macroblind}
\end{table}

We re-ran the full 5-window sensitivity with
\texttt{gate3\_mode=credit} (NFCI replacing MOVE; gate fires when
NFCI~$\geq$~0, i.e. any above-neutral tightness). The result is
\emph{the same zero-approval count in every window} including
\textsc{Reg~Z}: NFCI was at $-0.504$ on the peak-BSI day. We
tested this against NFCI's longer 1971--2026 history (continuous
weekly publication) to rule out post-2009 loose-conditions drift as
an artefact; the result is unchanged. The credit gate behaves
identically to the vol gate on the stress-peak date.

\paragraph{Architectural implication.}
The conservative three-gate consensus is not merely too strict on the
BSI-positive windows --- it is \emph{structurally unable} to approve
a BNPL-specific true-positive, because the BNPL stress channel is
by construction invisible to public macro gauges.
Two honest design responses remain:
\begin{enumerate}[leftmargin=*,itemsep=2pt]
\item \textbf{BSI-only high-conviction bypass.} If BSI~$z$ exceeds a
paper-specified super-threshold (e.g. $|z| > 10\sigma$, which has
triggered exactly once in eight years — on the Reg~Z deadline), admit
the trade without macro corroboration, with explicit Type-I risk
disclosure to the execution layer. This is the closest analog to the
way Paulson~\&~Co. traded the 2006--2007 subprime signal against a
macro tape that would not corroborate until months later.
\item \textbf{Degrade to a three-gate rule for BNPL-specific
events.} Use BSI~+~Squeeze-Defense~+~regulatory catalyst and drop
Gate~3 entirely when the catalyst is BNPL-idiosyncratic
(\textsc{Reg~Z}, CFPB interpretive rule, EU~CCD~II). This is a
taxonomy choice rather than a quantitative one.
\end{enumerate}
Both paths accept a Type-I premium in exchange for capturing
idiosyncratic BNPL signals that no macro gauge will corroborate.
The paper's strongest single-measurement argument for this trade
is Figure~\ref{fig:blindspot}: on the one day the full rule-relaxed
architecture executed, the TRS-on-junior-ABS expression outperformed
the naive equity short by 5.33 percentage points. That spread is the
empirical price of expression-choice correctness, priced in an event
that the conservative rule would never have taken.

Figure~\ref{fig:move} retains the long-horizon MOVE-vs-BSI overlay
for completeness. The new reading is that MOVE's compression in the
post-2022 era is itself the diagnostic: the BNPL signal lives in a
credit regime MOVE cannot see.

\begin{figure}[H]
\centering
\includegraphics[width=0.92\textwidth]{figures/fig6_move_vs_bsi.png}
\caption{MOVE index (crimson, right axis) versus BSI z-score (cyan,
left axis) over the full 2018--2026 sample. Dashed crimson line at
MOVE~=~120. The BSI spike on 17~January~2025 stands out against a
MOVE series that never sustainedly exceeds 120 in the post-2022
window; the divergence is the Subprime-2.0 blind-spot signature.}
\label{fig:move}
\end{figure}


%% =============================================================== Section 9
\section{Agentic Pod Architecture}
\label{sec:pod}

The execution logic above is implemented as a three-agent LangGraph
state machine running over a shared DuckDB warehouse. The three
agents are:

\begin{itemize}[leftmargin=2em]
\item \textbf{Macro Agent.} Reads the current BSI, MOVE, and
regulatory-catalyst calendar. Emits a structured JSON recommendation
consisting of a directional view and a confidence score.
\item \textbf{Quant Agent.} Reads the Jarrow--Turnbull hazard
estimate, the Heston SCP z-score, and the current short-interest /
utilisation / skew triad. Emits a second JSON recommendation with
an explicit TRS sizing suggestion conditional on the macro agent
approving.
\item \textbf{Risk Manager.} Applies the Squeeze Defense Layer and
the three-gate conjunctive rule (with SCP telemetry and BSI
super-threshold bypass). Writes the final trade decision to
\texttt{logs/agent\_decisions/} in JSONL and, via a unified LLM
client wrapper, composes a one-paragraph narrative explaining the
decision to a human trader.
\end{itemize}

The LLM client routes to NVIDIA NIM (Nemotron-4 8B for the Macro
Agent, Nemotron-Mini-4B-Instruct for the Risk Manager) as the primary
provider, with Anthropic Claude as a graceful-degradation fallback
for structured-output tasks. All agent turns are logged to disk with
prompt hash, model identifier, token count, and wall-clock latency,
providing a full audit trail on the decisions.

A React-based institutional terminal (served from a single static
\texttt{index.html}) subscribes to a \texttt{pod\_snapshot.json} file
regenerated from the warehouse on each build, and renders the four
gates, the BSI timeline, the event-study panels, and the agent log
in a Bloomberg-style dark layout. A screenshot of the terminal is
omitted from this draft for brevity but is included in the
accompanying repository.


%% =============================================================== Section 10
\section{Risk, Limitations, and Future Work}
\label{sec:risk}

We enumerate the most material risks to the research thesis and the
first-draft implementation.

\paragraph{Data limitations.}
Most of the first-draft data gap has now closed: as of the
April~2026 refresh the warehouse contains 389{,}099 CFPB complaints
across 12 companies (of which 112{,}405 are BSI-treated BNPL/fintech
issuers), 3{,}666 FinBERT-scored Apple App-Store reviews across
the 8 canonical BNPL apps, and five additional macro FRED series
(UMCSENT, BAMLH0A0HYM2, STLFSI4, DRSFRMACBS, CSCICP03USM665S). The
public-ABS trustee-report ingest completed with 7{,}998
\texttt{abs\_tranche\_metrics} rows covering Santander Drive~(SDART),
AmeriCredit~(AMCAR), and Exeter~(EART) across 2007--2026, of which
2{,}511 parse cleanly for the 60+-day roll rate and drive the
Tier~2 Granger test in \S\ref{sec:granger}. Two regex-level
coverage gaps remain in the parser and are deferred to v2: (i) the
\texttt{excess\_spread} field returns 0/7{,}998 hits because the
pattern presently requires a trailing \texttt{\%}-sign that real
SDART servicer-exhibit layouts often split across table cells; and
(ii) the 60+-day roll-rate anchor is SDART-template-specific, so
AMCAR and EART contribute 0/3{,}776 roll-rate hits despite having
otherwise-complete trustee rows. Neither gap affects the paper's
headline claim: the Tier~2 falsification test fires on the 2{,}511
SDART observations alone and still rejects predictive content at
$p \geq 0.95$. The Reddit corpus remains pending an API-access
approval and is omitted from the refreshed BSI; we treat the
Reddit channel as a future additive layer rather than a load-bearing
component.

\paragraph{Macro corroboration is structurally unavailable.}
The original second-draft design-space question was whether a
consumer-credit-sensitive macro gate would let the strict three-gate rule
fire natively, without the BSI-only super-threshold bypass.
\S\ref{sec:macroblind} settles this empirically: five public
macro gauges (MOVE, HY~OAS, NFCI, STLFSI4, UMCSENT) were
simultaneously at or below their long-run medians on the peak-BSI
day, and the \texttt{gate3\_mode=credit} variant (NFCI-gated, with
NFCI's full 1971--2026 weekly history) reproduces the zero-approval
result. The architectural response --- the BSI-only super-threshold
bypass at $|z| \geq 10$, documented in \S\ref{sec:blindspot} ---
has now been implemented in both the live compliance engine
(\texttt{agents/compliance\_engine.py}) and the backtest driver
(\texttt{backtest/event\_study.py::evaluate\_three\_gates}). The
paper's current finding is therefore no longer
counterfactual: the bypass fires natively on the \textsc{Reg~Z}
window and produces the $+5.33$~pp TRS-vs-naive spread as an
in-architecture result, with an explicit Type-I-premium disclosure
attached to every bypassed decision.

\paragraph{AFRMT trustee data is not public (and that is the thesis).}
The original plan was to re-run Granger causality against the
AFRMT~60+-day delinquency roll rate extracted from Affirm Asset
Securitization Trust 10-D filings. In implementing that pipeline we
discovered that Affirm's securitizations are \emph{144A private
placements}: Affirm Holdings (CIK 0001820953) has filed zero 10-D or
ABS-15G filings on EDGAR across the entire 2020--2026 sample. The
same holds for Afterpay (CIK 0001512673, Block, Inc. parent). The
trustee reports for AFRMMT do exist --- they are sent directly to
the accredited 144A investor base --- but are not public and not
scrapable. This is \emph{the thesis}: the BNPL-ABS secondary market
is definitionally opaque to the public, which is why a
BSI-vs-macro-gauge blind spot exists in the first place. We expand
the backtest infrastructure to cover public-ABS sponsors (SDART,
AMCAR, EART, CACC) as cross-sectional comparators and document the
AFRMMT gap as a structural feature of the Subprime-2.0 thesis
rather than an implementation gap.

\paragraph{Threshold calibration disclosure.}
The BSI $z$-threshold ($+1.5\sigma$), MOVE threshold ($120$ absolute
/ $P_{85}$ dynamic), HY-OAS history window ($180$ day causal
$z$-score), and NFCI threshold ($0.0$) are all \emph{ex-ante} defaults
chosen before running the 5-window sensitivity, not fit to realised
firing patterns. We include this disclosure because the paper reports
both positive and negative results on the same calibration: BSI fires
hard on \textsc{Reg~Z}, macro gauges uniformly do not, and the
zero-approval count is invariant across threshold specifications.
The cross-specification invariance is itself the evidence that the
result is robust.

\paragraph{Non-stationarity.}
Credit-cycle models are famously non-stationary. Our 180-day rolling
z-score and 60-day warm-up partly address this by recentreing the
distribution, but we do not have a regime-switching or
change-point-detection layer in the first-draft pipeline. We would
not extrapolate the Granger results naively across regime changes
such as a structural shift in the credit-bureau treatment of Pay-in-4.

\paragraph{Scraper fragility.}
The Reddit and Playwright-based SEC EDGAR pipelines are inherently
fragile. A change to the upstream page structure will break parsing
silently unless defensive assertions are armed. Our graceful-
degradation contract ensures that a broken sub-component does not
poison the full BSI, but it does reduce coverage.

\paragraph{Behavioural channel non-stationarity: the AI-search
dark-channel shift.}
The Google-Trends friction bucket is load-bearing for the
\emph{direct-admission} leg of the BSI
(\S\ref{sec:data}) --- queries of the form \emph{``i can't pay
affirm''}, \emph{``affirm hardship program''}, and the eight other
first-person distress phrases listed in the Trends taxonomy. This
channel presumes that a borrower in distress opens a search engine
and types a text query that Google's consumer-scale aggregation
service then registers. That presumption is weakening in real time.
Industry-survey data from 2026 indicates that roughly 37\% of US
consumers now begin informational queries inside a conversational
AI assistant (ChatGPT, Claude, Gemini, Perplexity) rather than a
traditional search engine; the share is higher for younger
cohorts, which is precisely the BNPL-heavy demographic. These
AI-mediated queries are \emph{not} surfaced in Google Trends or
any equivalent public aggregate: they land in private model-provider
logs that are not disclosed to third parties, rate-limited, or
sold as a commercial dataset. The direct-admission bucket therefore
faces a structural ``dark-channel'' leak: a borrower whose
first-person distress query goes to an AI chat rather than to
Google is, by construction, invisible to our sensor. This is a
harder version of the general-non-stationarity problem: the
\emph{channel itself} is migrating, not just the within-channel
distribution. Two mitigations are possible without closing the
gap fully. First, re-weighting the direct-admission bucket's
$\omega_{\text{DA}}$ coefficient downward as the AI-search share
rises, in proportion to a published industry-benchmark series,
would keep the BSI honest at the level of \emph{current} coverage
rather than \emph{historical} coverage. Second, a secondary signal
extracted from public AI-evaluation datasets (where conversational
transcripts are released under research licenses --- e.g.\ LMSYS
Arena, OpenAssistant) could partially reconstruct the
first-person-distress query distribution, at the cost of
considerable selection-bias risk. Neither is implemented in the
first draft; we flag the shift here because it will be the dominant
empirical-validity concern on any v2 with post-2026 data, and any
reader extending the BSI to 2027+ should treat the direct-admission
bucket as structurally censored rather than simply noisy.

\paragraph{Regulatory arbitrage.}
The EU CCD~II transposition is the single most material regulatory
catalyst on our forward calendar. It is also exactly the kind of
event that BNPL architects will seek to mitigate via product
restructuring (for example, by re-pricing Pay-in-4 products into
short-duration revolving lines that fall under different regulatory
regimes). A mitigated catalyst is a weaker one.

\paragraph{Strict-threshold Type-II risk.}
As discussed in \S\ref{sec:empirical}, the strict three-gate rule fires
zero days in our event study. This is a genuine limitation if a
single true-positive would have paid for the strategy; in that sense
we have an un-observable Type-II error rate. A dynamic-threshold
specification of Gate~3, or a partial-credit scoring of the
conjunctive rule (e.g.\ ``two of three gates plus a 0.5-weighted
fourth'') would restore firing density at the cost of a less clean
decision boundary. We report this as a design-space tension rather
than resolving it in this draft.

\paragraph{No execution integration.}
The first-draft pod does not include a broker integration. All P\&L
is simulated from historical public data. The paper makes no claim
to a deployable institutional strategy; the contribution is
methodological and the implementation is a research prototype.


%% =============================================================== Section 11
\section{Conclusion}
\label{sec:conclusion}

We have argued that BNPL credit, in its current regulatory and
reporting posture, exhibits structural features that recall the
opacity of early-2000s subprime mortgages --- bureau-invisibility,
accelerated ABS amortisation, off-ledger consumer stacking --- and
that the correct instrument for expressing a bearish view is a
Total-Return Swap on a junior BNPL ABS tranche rather than a
directional short of the architect-tier equity complex. We have
built a reproducible daily BNPL Stress Index from public and
alternative data and tested whether it Granger-predicts HYG
high-yield credit stress at 4--8 week horizons. We
\emph{failed to reject the null} ($p > 0.69$ at every tested lag).
We argue throughout the paper that this null is the
thesis-consistent outcome: if a BNPL-specific psychological
stress signal were Granger-predictive of broad-market high-yield
credit at monthly horizons, the Subprime-2.0 opacity claim would
be falsified. An earlier specification in which the composite was
dominated by the MOVE sub-component did reject the null at
$p < 10^{-4}$ and was reported as such in a prior draft; we now
regard that result as a MOVE-artefact and retract it
(see \S\ref{sec:granger}). The honest empirical contribution is
therefore (i) the composite itself, (ii) the 5-macro-gauge
blind-spot signature at the Reg~Z peak, and (iii) the implemented
BSI-only super-threshold bypass that lets the system approve the
single largest stress event in the sample
(\S\ref{sec:blindspot}). We implement all of this inside a
three-gate (BSI $\times$ MOVE $\times$ CCD~II) execution framework
with the equity-derived Squeeze-Compression Premium demoted to
non-gating telemetry --- since the expression is TRS on a junior
ABS tranche rather than an AFRM equity short, there is no squeeze
leg for an SCP gate to protect. Our empirical event-study over
five 2022--2025 episodes finds zero approved days under the strict
three-gate rule and three approved days under the $|\bsi_z|\geq10$
bypass (all during the Reg~Z cluster), which we interpret as
correct Type-II conservatism on macro-visible stress and an
explicit, audited Type-I premium for the one episode where
BNPL-specific stress exceeded every macro gauge's resolution.

The most direct extensions of this work are (i) backfilling the
BSI sub-components that are implemented but not yet ingested at
scale, (ii) re-running the Granger tests against the
theoretically-preferred AFRMT roll-rate target once trustee data
are available, and (iii) characterising the sensitivity of the
three-gate firing rate to dynamic-threshold specifications of the
MOVE gate.

Beyond those immediate extensions, the most interesting open
questions are structural: whether the BNPL credit cycle is
genuinely orthogonal to the revolving-card cycle, or whether the
two merely appear orthogonal because they are observed through
different instruments. If BNPL stress turns out to lead rather
than co-move with revolving-card stress, the research agenda is
far larger than this paper; if the two are in fact the same latent
process sampled differently, our results become a statement about
measurement technology rather than about a new credit regime. That
question cannot be answered from the 2018--2026 sample alone, but
the methodological stack we have built is the right one to answer
it.


%% =============================================================== References
\bibliography{references}


%% =============================================================== Appendix A
\appendix
\section{Reproducibility}
\label{app:repro}

The full codebase is a Python package with a single-file
\texttt{Makefile} that chains all stages:

\begin{itemize}[leftmargin=2em,noitemsep]
\item \texttt{make setup} --- creates a \texttt{uv} virtual
environment and installs the package in editable mode with
development extras.
\item \texttt{make ingest} --- runs every data-ingestion module
(FRED, SEC EDGAR, CFPB, Reddit, Google Trends, options chain,
short interest) into DuckDB.
\item \texttt{make bsi} --- computes the daily BSI and persists
it to \texttt{bsi\_daily}.
\item \texttt{make validate} --- runs Granger causality and
rolling out-of-sample tests.
\item \texttt{make backtest} --- runs the event-study simulator
and emits per-window P\&L CSVs.
\item \texttt{make pod} --- executes one end-to-end pass of the
LangGraph pod, logging every agent decision.
\item \texttt{make snapshot} --- rebuilds the React terminal's
\texttt{pod\_snapshot.json}.
\item \texttt{make paper} --- regenerates every figure in this
document and compiles the LaTeX.
\end{itemize}

Each target is idempotent and writes intermediate CSVs / Parquet
to disk, so downstream stages can be re-run against cached
upstream outputs. A clean clone-and-run on a machine with Python
3.11+, \texttt{uv}, and a working \TeX{} distribution should
reproduce every number and figure in this paper via:
\begin{center}
\texttt{git clone <repo> \&\& cd bnpl-pod \&\& make setup ingest bsi validate backtest paper}
\end{center}

\end{document}
